\input{../tex-inputs/latex-leseansicht-vorspann}

\section[Gustav Schwarzkopf an Arthur Schnitzler, 28. 7. 1901]{L04305 Gustav Schwarzkopf an Arthur Schnitzler, 28. 7. 1901}
\nopagebreak\mylabel{L04305v}
\rehead{ }\normalsize\beginnumbering\briefempfaengerindex{Schnitzler, Arthur@\textsc{Schnitzler, Arthur}!zzzSchwarzkopf, Gustav@\emph{von Gustav Schwarzkopf}!1901-07-282@{28. 7. 1901}|(be}
\toendnotes[C]{\smallbreak\pagebreak[2]}
\correspDesc{Versand  durch Gustav Schwarzkopf am 28. 7. 1901 in Wien
\newline{}Übermittlung  am 29. 7. 1901 in Wien
\newline{}Zustellung  am 30. 7. 1901 in Vahrn
\newline{}Erhalt  durch Arthur Schnitzler am 30. 7. 1901 in Vahrn}\toendnotes[C]{\smallbreak}
\Standort{CUL, Schnitzler, B 96a.}
\physDesc{Postkarte, 563 Zeichen
\newline{}Handschrift: schwarze Tinte, deutsche Kurrent
\newline{}Versand: 1) Stempel: »\nobreak{}\oindex{I., Innere Stadt@\textbf{I., Innere Stadt}, \emph{Verwaltungsgebiet}|pwk}Wien 1/1 1, 29. 7. 01, 8–9V\nobreak{}«.   2) Stempel: »\nobreak{}\oindex{Vahrn@\textbf{Vahrn}, \emph{Hauptstadt}|pwk}Vahrn, 30. 7. {[}01{]}\nobreak{}«. }\toendnotes[C]{\smallbreak}\pstart{}{\pb}Herrn \textsc{D\textsuperscript{r} Arthur Schnitzler}\pend{}\pstart{}\textsc{Vahren} bei \textsc{Brixen}\oindex{Vahrn@\textbf{Vahrn}, \emph{Hauptstadt}|pw}\textsc{Brixen}\pend{}\pstart{}\textsc{Pension Waldsacker\oindex{Gasthof und Pension Waldsacker@\textbf{Gasthof und Pension Waldsacker}, \emph{Beherbergungsgebäude}|pw}}\pend{}{\bigskip}\vspace{1em}
\pstart
           \raggedleft{}{\pb}Wien\oindex{Wien@\textbf{Wien}, \emph{Verwaltungsgebiet}|pw}{ }28. VII 01\pend
           \vspace{0.5em}
\pstart
           Lieber Arthur, ich habe Ihnen allerdings gar nichts mitzuteilen,
               will Ihnen nur für Ihren \label{K_L04305-1v}\edtext{Brief}{\lemma{\textnormal{\emph{Brief}}}\Cendnote{\textnormal{XXXX Auszeichnungsfehler: Dokument L04215 nicht gefunden.
               }}}\label{K_L04305-1} danken und ihre Grüße erwidern. Mein Bruder\pwindex{Schwarzkopf, Max 12.\,6.\,1857 Wien – 14.\,4.\,1928 ebd.@\textsc{Schwarzkopf, Max} (12.\,6.\,1857 Wien – 14.\,4.\,1928 ebd.), \emph{Rechtsanwalt}|pwv}, der erſt in den erſten Auguſttagen nach Wien\oindex{Wien@\textbf{Wien}, \emph{Verwaltungsgebiet}|pw} ko{\geminationm}t, wird Ihnen gewiß gerne{ }ſeine juridischen Ke{\geminationn}tniſſe zu Dienſten{ }ſtellen. –
               Vorgeſtern war ich Vormittag im Prater\oindex{Wien@\textbf{Wien}!II., Leopoldstadt@\textbf{II., Leopoldstadt}!Prater@\textbf{Prater}, \emph{Park}|pw} mit — Goldma{\geminationn}\pwindex{Goldmann, Paul 31.\,1.\,1865 Breslau – 25.\,9.\,1935 Wien@\textsc{Goldmann, Paul} (31.\,1.\,1865 Breslau – 25.\,9.\,1935 Wien), \emph{Schriftsteller, Journalist}|pw}, den ich zufällig getroffen habe. Ich war auch
               bei der \label{K_L04305-2v}\edtext{Trauung der{ }ſchönen \textsc{Hetsey\pwindex{Hétsey-Holzer, Alice 3.\,9.\,1875 Wien – 11.\,2.\,1939 ebd.@\textsc{Hétsey-Holzer, Alice} (3.\,9.\,1875 Wien – 11.\,2.\,1939 ebd.), \emph{Schauspielerin}|pw}}}{\lemma{\textnormal{\emph{Trauung … Hetsey}}}\Cendnote{\textnormal{Alice Hétsey\pwindex{Hétsey-Holzer, Alice 3.\,9.\,1875 Wien – 11.\,2.\,1939 ebd.@\textsc{Hétsey-Holzer, Alice} (3.\,9.\,1875 Wien – 11.\,2.\,1939 ebd.), \emph{Schauspielerin}|pwk} und 
                  Rudolf Holzer\pwindex{Holzer, Rudolf 28.\,7.\,1875 Wien – 17.\,7.\,1965 ebd.@\textsc{Holzer, Rudolf} (28.\,7.\,1875 Wien – 17.\,7.\,1965 ebd.), \emph{Schriftsteller, Journalist}|pwk} heirateten am 15. 7. 1901
                  in Wien\oindex{Wien@\textbf{Wien}, \emph{Verwaltungsgebiet}|pwk}.
               }}}\label{K_L04305-2} als Repräſenheit von »ganz Wien\oindex{Wien@\textbf{Wien}, \emph{Verwaltungsgebiet}|pw}«?\pend
           
\pstart
           Empfehlen Sie mich beſtens den beiden Damen\pwindex{Schnitzler, Olga 17.\,1.\,1882 Wien – 13.\,1.\,1970 Lugano@\textsc{Schnitzler, Olga} (17.\,1.\,1882 Wien – 13.\,1.\,1970 Lugano), \emph{Schauspielerin, Sängerin}|pwv}\pwindex{Steinrück, Elisabeth 19.\,11.\,1885 – 7.\,4.\,1920 Partenkirchen@\textsc{Steinrück, Elisabeth} (19.\,11.\,1885 – 7.\,4.\,1920 Partenkirchen)|pwv}.\pend
           
\pstart
           Herzlichſt Ihr{\\[\baselineskip]}\spacefill\mbox{Gustav}\pend
           \leftskip=0em{}\selectlanguage{ngerman}\endnumbering\briefempfaengerindex{Schnitzler, Arthur@\textsc{Schnitzler, Arthur}!zzzSchwarzkopf, Gustav@\emph{von Gustav Schwarzkopf}!1901-07-282@{28. 7. 1901}|)be}\mylabel{L04305h}
\begin{anhang}
\end{anhang}\newcommand{\dateiname}{L04305}\newcommand{\titel}{Gustav Schwarzkopf an Arthur Schnitzler, 28. 7. 1901}\newcommand{\editorInnen}{Herausgegeben von Jahnke, SelmaMüller, Martin Anton}\input{../tex-inputs/latex-leseansicht-abspann}
